\section{结论}

本文围绕视觉刺激条件下的脑电信号分析与计算建模，构建从数据预处理到神经机制解释的分层模型框架。

针对问题一，建立了保真约束的脑电预处理与ERP提取方法。结果表明，质量加权能够提升共同ERP的稳定性，而更复杂的学习型平滑并未在所有人工信号保真测试中通过验收，因此最终采用基础保真处理作为后续统一输入。这一结果说明，在弱脑电特征分析中，模型复杂度并非越高越好，视觉任务相关信息的保留应优先于单纯的平滑程度。

\TODO{问题二完成后补写LGN--皮层--头皮生成机制的主要结论。}

\TODO{问题三完成后补写认知模型的主要结论与脑机接口/精神性疾病研究意义。}
